%% kapa: exam=1 round=35 year=2024 period=1 form=A subject=real_estate subject_ko=부동산학원론 seq=2 number=82
%% source: Q-Net 감정평가사 기출문제 (cst00309) artlSeq=5213229 file=raw/1008/5213229/1교시 문제지 원본.pdf
%% conversion: kordoc 4.24.51 -> tools/gen1.py (auto); review=proofread-sonnet; flags=none
\begin{question}{82}
공인중개사법령상 개업공인중개사에 관한 내용으로 옳지 않은 것은?
\begin{choices}
  \item 개업공인중개사는 그 사무소의 명칭에 “공인중개사사무소” 또는 “부동산중개”라는 문자를 사용하여야 한다.
  \item 개업공인중개사가 아닌 자는 중개대상물에 대한 표시ㆍ광고를 하여서는 아니 된다.
  \item 개업공인중개사는 「민사집행법」에 의한 경매 및 「국세징수법」 그 밖의 법령에 의한 공매대상 부동산에 대한 권리분석 및 취득의 알선과 매수신청 또는 입찰신청의 대리를 할 수 있다.
  \item 개업공인중개사는 대통령령으로 정하는 기준과 절차에 따라 등록관청의 허가를 받아 그 관할 구역 외의 지역에 분사무소를 둘 수 있다.
  \item 개업공인중개사는 다른 사람에게 자기의 성명 또는 상호를 사용하여 중개업무를 하게 하거나 자기의 중개사무소등록증을 양도 또는 대여하는 행위를 하여서는 아니된다.
\end{choices}
\end{question}
\answer{4}
